% =====================================================================
% m3.tex -- Member 3 (Main-Memory Interface & Top-Level Integration) report sections
%   * Literature Review, subsection C (write policies and memory interfaces)
%   * Proposed Solution, subsection C (memory interface, refill & top-level integration)
%
% For M5: % =====================================================================
% m3.tex -- Member 3 (Main-Memory Interface & Top-Level Integration) report sections
%   * Literature Review, subsection C (write policies and memory interfaces)
%   * Proposed Solution, subsection C (memory interface, refill & top-level integration)
%
% For M5: % =====================================================================
% m3.tex -- Member 3 (Main-Memory Interface & Top-Level Integration) report sections
%   * Literature Review, subsection C (write policies and memory interfaces)
%   * Proposed Solution, subsection C (memory interface, refill & top-level integration)
%
% For M5: % =====================================================================
% m3.tex -- Member 3 (Main-Memory Interface & Top-Level Integration) report sections
%   * Literature Review, subsection C (write policies and memory interfaces)
%   * Proposed Solution, subsection C (memory interface, refill & top-level integration)
%
% For M5: \input{sections/m3} pieces into the IEEEtran report, or copy the
% blocks marked BEGIN/END. Citation keys match m1.tex.
% =====================================================================

% ---------------------- BEGIN LITERATURE REVIEW C ---------------------
\subsection{Write Policies and Main-Memory Interfaces}
\label{sec:litreview-m3}

The interaction between a cache and main memory is governed by two
orthogonal policy choices: \emph{write policy} (when memory is updated) and
\emph{write-allocation policy} (whether a line is brought into cache on a
write miss)~\cite{smith1982, hennessy2017}.

\subsubsection{Write-Through vs.\ Write-Back}
Under a \textbf{write-through} policy, every CPU write immediately updates
both the cache line (if present) and the backing store. The primary
advantage is architectural simplicity and data consistency: main memory is
always up to date, eliminating the need for dirty bits or writebacks on
evictions. Consequently, handling a read-miss refill simply overwrites the
selected line without an eviction writeback~\cite{jouppi1990}. The drawback
is elevated memory-bus traffic, as every store instruction generates a memory
bus transaction regardless of temporal write locality~\cite{patterson2020}.

In contrast, a \textbf{write-back} policy defers memory updates until a
dirty line is evicted from the cache. This significantly reduces write
traffic—particularly for tight inner loops repeatedly modifying local
variables—at the expense of maintaining a per-line \emph{dirty bit} and
implementing a two-phase eviction-refill controller. Phase~II of our project
extends this direct-mapped controller to support write-back.

\subsubsection{Write-Allocate vs.\ No-Write-Allocate}
On a \textbf{write miss}, a \emph{write-allocate} (fetch-on-write) policy
fetches the target block from memory into cache before completing the write,
exploiting spatial and temporal locality under the assumption of subsequent
accesses. Conversely, \emph{no-write-allocate} (write-around) forwards the
write directly to memory, leaving the cache unchanged~\cite{hennessy2017, hill1989}.
In Phase~I, we combine write-through with no-write-allocate: the cleanest and
most deterministic pairing, as writes never trigger block refills or dirty-bit
management.

\subsubsection{Memory Interface Protocols and Burst Latencies}
Modern synchronous DRAM architectures access memory in \emph{bursts}: an
initial row-activation and column-selection latency ($t_{\mathrm{RCD}} + \text{CL}$)
is paid for the first word, after which consecutive words within the open row
are transferred at one word per clock cycle~\cite{hennessy2017}. Our behavioural
memory model parameterises this asymmetry via \texttt{MEM\_FIRST\_LAT} (10 cycles)
and \texttt{MEM\_NEXT\_LAT} (1 cycle). A 4-word (16-byte) block refill thus requires
$10 + 3 \times 1 = 13$ memory cycles, forming the baseline miss penalty for the
AMAT calculations detailed in Section~\ref{sec:analysis-m5}.
% ----------------------- END LITERATURE REVIEW C ----------------------


% --------------------- BEGIN PROPOSED SOLUTION C ----------------------
\subsection{Memory Interface Bridge and Top-Level Integration}
\label{sec:proposed-m3}

\subsubsection{Main-Memory Behavioural Model (\texttt{main\_memory.v})}
The main memory is implemented as a parameterised behavioural model with a
flat 64\,K-word (256\,KB) register array pre-initialised with predictable
word-aligned data ($M[k] = 4k$). The module implements a 3-state controller
(\texttt{IDLE} $\to$ \texttt{FIRST} $\to$ \texttt{BURST}):
\begin{itemize}
    \item \textbf{Block Read:} On \texttt{mem\_req} with \texttt{mem\_wr = 0},
        it latches the block-aligned address and counts \texttt{MEM\_FIRST\_LAT}
        cycles for word~0, followed by \texttt{MEM\_NEXT\_LAT} cycles each for
        words 1 through 3. When all 4 words are assembled, \texttt{mem_ready}
        pulses for 1 cycle with the full 128-bit block on \texttt{mem\_rdata}.
    \item \textbf{Single-Word Write:} On \texttt{mem\_wr = 1}, it waits
        \texttt{MEM\_FIRST\_LAT} cycles, commits the 32-bit word to the array,
        and pulses \texttt{mem\_ready}.
\end{itemize}

\subsubsection{Memory Interface Bridge (\texttt{mem\_interface.v})}
The memory interface decouples the controller FSM from main memory timing. It
latches outgoing requests at the \texttt{IDLE} $\to$ \texttt{BUSY} transition,
asserts \texttt{mm\_req} for one cycle to initiate memory access, holds address
and write-data stable, and routes the returned 128-bit line combinationally
to \texttt{data\_ram.wdata\_line} via \texttt{refill\_data}.

Crucially, \texttt{fsm\_mem\_ready} is gated by the interface state
($\texttt{mm\_ready} \wedge \texttt{mi\_state == BUSY}$). Because \texttt{mm\_req}
is de-asserted once accepted, spurious re-triggering is completely prevented
when the FSM transitions from \texttt{ALLOCATE} to \texttt{UPDATE}.

\subsubsection{Top-Level Integration (\texttt{cache\_top.v})}
\texttt{cache\_top} connects all eight sub-modules into a cohesive system,
exposing only CPU-side signals:
\begin{equation}
    \{\texttt{clk}, \texttt{rst}, \texttt{cpu\_req}, \texttt{cpu\_wr}, \texttt{cpu\_addr}, \texttt{cpu\_wdata}\}
    \longrightarrow
    \{\texttt{cpu\_rdata}, \texttt{cpu\_ready}\}
\end{equation}

The cycle-by-cycle operational flows are verified as follows:
\begin{enumerate}
    \item \textbf{Read Hit (1 cycle):} With combinational tag and valid lookup,
        \texttt{hit} asserts in \texttt{IDLE}; \texttt{cpu\_ready} fires in the
        same cycle (fast path).
    \item \textbf{Read Miss (19 cycles total):} \texttt{IDLE} (1) $\to$
        \texttt{COMPARE\_TAG} (1) $\to$ \texttt{ALLOCATE} (14 cycles: 1 interface setup
        + 13 memory burst) $\to$ \texttt{UPDATE} (1 cycle to write data line, tag,
        and valid bit) $\to$ \texttt{RESPOND} (1 cycle asserting \texttt{cpu\_ready})
        $\to$ \texttt{IDLE}.
    \item \textbf{Write Hit (1 cycle):} Fast-path update of the cache word in
        \texttt{data\_ram} while concurrently dispatching the write to memory.
    \item \textbf{Write Miss (16 cycles total):} \texttt{COMPARE\_TAG} $\to$
        \texttt{ALLOCATE} (10 cycles write to memory) $\to$ \texttt{RESPOND}. The
        cache tag and data arrays remain untouched (no-write-allocate).
\end{enumerate}

\subsubsection{Unit and Integration Verification}
Member~3 deliverables are verified via an automated test suite (\texttt{tb/run\_m3\_tests.sh}):
\begin{itemize}
    \item \texttt{tb\_main\_memory.v}: Validates exact 13-cycle read burst,
        10-cycle write, and data integrity across all words.
    \item \texttt{tb\_mem\_interface.v}: Validates request latching, ready pulse
        generation, and absence of phantom re-triggering.
    \item \texttt{tb\_cache\_top.v}: Validates full integration over cold misses,
        spatial locality hits, write-through, no-write-allocate, and conflict-miss
        thrashing (addresses 1\,KB apart).
\end{itemize}
% ---------------------- END PROPOSED SOLUTION C -----------------------


% -------------- REFERENCES (merge into the shared list) ---------------
% \begin{thebibliography}{9}
% \bibitem{smith1982} A.~J. Smith, ``Cache memories,'' \emph{ACM Computing
%   Surveys}, vol.~14, no.~3, pp.~473--530, Sep. 1982.
% \bibitem{hill1989} M.~D. Hill and A.~J. Smith, ``Evaluating associativity in
%   CPU caches,'' \emph{IEEE Trans. Computers}, vol.~38, no.~12,
%   pp.~1612--1630, Dec. 1989.
% \bibitem{jouppi1990} N.~P. Jouppi, ``Improving direct-mapped cache
%   performance by the addition of a small fully-associative cache and
%   prefetch buffers,'' in \emph{Proc. 17th Annu. Int. Symp. Computer
%   Architecture (ISCA)}, 1990, pp.~364--373.
% \bibitem{hennessy2017} J.~L. Hennessy and D.~A. Patterson, \emph{Computer
%   Architecture: A Quantitative Approach}, 6th~ed. Cambridge, MA, USA:
%   Morgan Kaufmann, 2017.
% \bibitem{patterson2020} D.~A. Patterson and J.~L. Hennessy, \emph{Computer
%   Organization and Design RISC-V Edition}, 2nd~ed. Morgan Kaufmann, 2020.
% \end{thebibliography}
 pieces into the IEEEtran report, or copy the
% blocks marked BEGIN/END. Citation keys match m1.tex.
% =====================================================================

% ---------------------- BEGIN LITERATURE REVIEW C ---------------------
\subsection{Write Policies and Main-Memory Interfaces}
\label{sec:litreview-m3}

The interaction between a cache and main memory is governed by two
orthogonal policy choices: \emph{write policy} (when memory is updated) and
\emph{write-allocation policy} (whether a line is brought into cache on a
write miss)~\cite{smith1982, hennessy2017}.

\subsubsection{Write-Through vs.\ Write-Back}
Under a \textbf{write-through} policy, every CPU write immediately updates
both the cache line (if present) and the backing store. The primary
advantage is architectural simplicity and data consistency: main memory is
always up to date, eliminating the need for dirty bits or writebacks on
evictions. Consequently, handling a read-miss refill simply overwrites the
selected line without an eviction writeback~\cite{jouppi1990}. The drawback
is elevated memory-bus traffic, as every store instruction generates a memory
bus transaction regardless of temporal write locality~\cite{patterson2020}.

In contrast, a \textbf{write-back} policy defers memory updates until a
dirty line is evicted from the cache. This significantly reduces write
traffic—particularly for tight inner loops repeatedly modifying local
variables—at the expense of maintaining a per-line \emph{dirty bit} and
implementing a two-phase eviction-refill controller. Phase~II of our project
extends this direct-mapped controller to support write-back.

\subsubsection{Write-Allocate vs.\ No-Write-Allocate}
On a \textbf{write miss}, a \emph{write-allocate} (fetch-on-write) policy
fetches the target block from memory into cache before completing the write,
exploiting spatial and temporal locality under the assumption of subsequent
accesses. Conversely, \emph{no-write-allocate} (write-around) forwards the
write directly to memory, leaving the cache unchanged~\cite{hennessy2017, hill1989}.
In Phase~I, we combine write-through with no-write-allocate: the cleanest and
most deterministic pairing, as writes never trigger block refills or dirty-bit
management.

\subsubsection{Memory Interface Protocols and Burst Latencies}
Modern synchronous DRAM architectures access memory in \emph{bursts}: an
initial row-activation and column-selection latency ($t_{\mathrm{RCD}} + \text{CL}$)
is paid for the first word, after which consecutive words within the open row
are transferred at one word per clock cycle~\cite{hennessy2017}. Our behavioural
memory model parameterises this asymmetry via \texttt{MEM\_FIRST\_LAT} (10 cycles)
and \texttt{MEM\_NEXT\_LAT} (1 cycle). A 4-word (16-byte) block refill thus requires
$10 + 3 \times 1 = 13$ memory cycles, forming the baseline miss penalty for the
AMAT calculations detailed in Section~\ref{sec:analysis-m5}.
% ----------------------- END LITERATURE REVIEW C ----------------------


% --------------------- BEGIN PROPOSED SOLUTION C ----------------------
\subsection{Memory Interface Bridge and Top-Level Integration}
\label{sec:proposed-m3}

\subsubsection{Main-Memory Behavioural Model (\texttt{main\_memory.v})}
The main memory is implemented as a parameterised behavioural model with a
flat 64\,K-word (256\,KB) register array pre-initialised with predictable
word-aligned data ($M[k] = 4k$). The module implements a 3-state controller
(\texttt{IDLE} $\to$ \texttt{FIRST} $\to$ \texttt{BURST}):
\begin{itemize}
    \item \textbf{Block Read:} On \texttt{mem\_req} with \texttt{mem\_wr = 0},
        it latches the block-aligned address and counts \texttt{MEM\_FIRST\_LAT}
        cycles for word~0, followed by \texttt{MEM\_NEXT\_LAT} cycles each for
        words 1 through 3. When all 4 words are assembled, \texttt{mem_ready}
        pulses for 1 cycle with the full 128-bit block on \texttt{mem\_rdata}.
    \item \textbf{Single-Word Write:} On \texttt{mem\_wr = 1}, it waits
        \texttt{MEM\_FIRST\_LAT} cycles, commits the 32-bit word to the array,
        and pulses \texttt{mem\_ready}.
\end{itemize}

\subsubsection{Memory Interface Bridge (\texttt{mem\_interface.v})}
The memory interface decouples the controller FSM from main memory timing. It
latches outgoing requests at the \texttt{IDLE} $\to$ \texttt{BUSY} transition,
asserts \texttt{mm\_req} for one cycle to initiate memory access, holds address
and write-data stable, and routes the returned 128-bit line combinationally
to \texttt{data\_ram.wdata\_line} via \texttt{refill\_data}.

Crucially, \texttt{fsm\_mem\_ready} is gated by the interface state
($\texttt{mm\_ready} \wedge \texttt{mi\_state == BUSY}$). Because \texttt{mm\_req}
is de-asserted once accepted, spurious re-triggering is completely prevented
when the FSM transitions from \texttt{ALLOCATE} to \texttt{UPDATE}.

\subsubsection{Top-Level Integration (\texttt{cache\_top.v})}
\texttt{cache\_top} connects all eight sub-modules into a cohesive system,
exposing only CPU-side signals:
\begin{equation}
    \{\texttt{clk}, \texttt{rst}, \texttt{cpu\_req}, \texttt{cpu\_wr}, \texttt{cpu\_addr}, \texttt{cpu\_wdata}\}
    \longrightarrow
    \{\texttt{cpu\_rdata}, \texttt{cpu\_ready}\}
\end{equation}

The cycle-by-cycle operational flows are verified as follows:
\begin{enumerate}
    \item \textbf{Read Hit (1 cycle):} With combinational tag and valid lookup,
        \texttt{hit} asserts in \texttt{IDLE}; \texttt{cpu\_ready} fires in the
        same cycle (fast path).
    \item \textbf{Read Miss (19 cycles total):} \texttt{IDLE} (1) $\to$
        \texttt{COMPARE\_TAG} (1) $\to$ \texttt{ALLOCATE} (14 cycles: 1 interface setup
        + 13 memory burst) $\to$ \texttt{UPDATE} (1 cycle to write data line, tag,
        and valid bit) $\to$ \texttt{RESPOND} (1 cycle asserting \texttt{cpu\_ready})
        $\to$ \texttt{IDLE}.
    \item \textbf{Write Hit (1 cycle):} Fast-path update of the cache word in
        \texttt{data\_ram} while concurrently dispatching the write to memory.
    \item \textbf{Write Miss (16 cycles total):} \texttt{COMPARE\_TAG} $\to$
        \texttt{ALLOCATE} (10 cycles write to memory) $\to$ \texttt{RESPOND}. The
        cache tag and data arrays remain untouched (no-write-allocate).
\end{enumerate}

\subsubsection{Unit and Integration Verification}
Member~3 deliverables are verified via an automated test suite (\texttt{tb/run\_m3\_tests.sh}):
\begin{itemize}
    \item \texttt{tb\_main\_memory.v}: Validates exact 13-cycle read burst,
        10-cycle write, and data integrity across all words.
    \item \texttt{tb\_mem\_interface.v}: Validates request latching, ready pulse
        generation, and absence of phantom re-triggering.
    \item \texttt{tb\_cache\_top.v}: Validates full integration over cold misses,
        spatial locality hits, write-through, no-write-allocate, and conflict-miss
        thrashing (addresses 1\,KB apart).
\end{itemize}
% ---------------------- END PROPOSED SOLUTION C -----------------------


% -------------- REFERENCES (merge into the shared list) ---------------
% \begin{thebibliography}{9}
% \bibitem{smith1982} A.~J. Smith, ``Cache memories,'' \emph{ACM Computing
%   Surveys}, vol.~14, no.~3, pp.~473--530, Sep. 1982.
% \bibitem{hill1989} M.~D. Hill and A.~J. Smith, ``Evaluating associativity in
%   CPU caches,'' \emph{IEEE Trans. Computers}, vol.~38, no.~12,
%   pp.~1612--1630, Dec. 1989.
% \bibitem{jouppi1990} N.~P. Jouppi, ``Improving direct-mapped cache
%   performance by the addition of a small fully-associative cache and
%   prefetch buffers,'' in \emph{Proc. 17th Annu. Int. Symp. Computer
%   Architecture (ISCA)}, 1990, pp.~364--373.
% \bibitem{hennessy2017} J.~L. Hennessy and D.~A. Patterson, \emph{Computer
%   Architecture: A Quantitative Approach}, 6th~ed. Cambridge, MA, USA:
%   Morgan Kaufmann, 2017.
% \bibitem{patterson2020} D.~A. Patterson and J.~L. Hennessy, \emph{Computer
%   Organization and Design RISC-V Edition}, 2nd~ed. Morgan Kaufmann, 2020.
% \end{thebibliography}
 pieces into the IEEEtran report, or copy the
% blocks marked BEGIN/END. Citation keys match m1.tex.
% =====================================================================

% ---------------------- BEGIN LITERATURE REVIEW C ---------------------
\subsection{Write Policies and Main-Memory Interfaces}
\label{sec:litreview-m3}

The interaction between a cache and main memory is governed by two
orthogonal policy choices: \emph{write policy} (when memory is updated) and
\emph{write-allocation policy} (whether a line is brought into cache on a
write miss)~\cite{smith1982, hennessy2017}.

\subsubsection{Write-Through vs.\ Write-Back}
Under a \textbf{write-through} policy, every CPU write immediately updates
both the cache line (if present) and the backing store. The primary
advantage is architectural simplicity and data consistency: main memory is
always up to date, eliminating the need for dirty bits or writebacks on
evictions. Consequently, handling a read-miss refill simply overwrites the
selected line without an eviction writeback~\cite{jouppi1990}. The drawback
is elevated memory-bus traffic, as every store instruction generates a memory
bus transaction regardless of temporal write locality~\cite{patterson2020}.

In contrast, a \textbf{write-back} policy defers memory updates until a
dirty line is evicted from the cache. This significantly reduces write
traffic—particularly for tight inner loops repeatedly modifying local
variables—at the expense of maintaining a per-line \emph{dirty bit} and
implementing a two-phase eviction-refill controller. Phase~II of our project
extends this direct-mapped controller to support write-back.

\subsubsection{Write-Allocate vs.\ No-Write-Allocate}
On a \textbf{write miss}, a \emph{write-allocate} (fetch-on-write) policy
fetches the target block from memory into cache before completing the write,
exploiting spatial and temporal locality under the assumption of subsequent
accesses. Conversely, \emph{no-write-allocate} (write-around) forwards the
write directly to memory, leaving the cache unchanged~\cite{hennessy2017, hill1989}.
In Phase~I, we combine write-through with no-write-allocate: the cleanest and
most deterministic pairing, as writes never trigger block refills or dirty-bit
management.

\subsubsection{Memory Interface Protocols and Burst Latencies}
Modern synchronous DRAM architectures access memory in \emph{bursts}: an
initial row-activation and column-selection latency ($t_{\mathrm{RCD}} + \text{CL}$)
is paid for the first word, after which consecutive words within the open row
are transferred at one word per clock cycle~\cite{hennessy2017}. Our behavioural
memory model parameterises this asymmetry via \texttt{MEM\_FIRST\_LAT} (10 cycles)
and \texttt{MEM\_NEXT\_LAT} (1 cycle). A 4-word (16-byte) block refill thus requires
$10 + 3 \times 1 = 13$ memory cycles, forming the baseline miss penalty for the
AMAT calculations detailed in Section~\ref{sec:analysis-m5}.
% ----------------------- END LITERATURE REVIEW C ----------------------


% --------------------- BEGIN PROPOSED SOLUTION C ----------------------
\subsection{Memory Interface Bridge and Top-Level Integration}
\label{sec:proposed-m3}

\subsubsection{Main-Memory Behavioural Model (\texttt{main\_memory.v})}
The main memory is implemented as a parameterised behavioural model with a
flat 64\,K-word (256\,KB) register array pre-initialised with predictable
word-aligned data ($M[k] = 4k$). The module implements a 3-state controller
(\texttt{IDLE} $\to$ \texttt{FIRST} $\to$ \texttt{BURST}):
\begin{itemize}
    \item \textbf{Block Read:} On \texttt{mem\_req} with \texttt{mem\_wr = 0},
        it latches the block-aligned address and counts \texttt{MEM\_FIRST\_LAT}
        cycles for word~0, followed by \texttt{MEM\_NEXT\_LAT} cycles each for
        words 1 through 3. When all 4 words are assembled, \texttt{mem_ready}
        pulses for 1 cycle with the full 128-bit block on \texttt{mem\_rdata}.
    \item \textbf{Single-Word Write:} On \texttt{mem\_wr = 1}, it waits
        \texttt{MEM\_FIRST\_LAT} cycles, commits the 32-bit word to the array,
        and pulses \texttt{mem\_ready}.
\end{itemize}

\subsubsection{Memory Interface Bridge (\texttt{mem\_interface.v})}
The memory interface decouples the controller FSM from main memory timing. It
latches outgoing requests at the \texttt{IDLE} $\to$ \texttt{BUSY} transition,
asserts \texttt{mm\_req} for one cycle to initiate memory access, holds address
and write-data stable, and routes the returned 128-bit line combinationally
to \texttt{data\_ram.wdata\_line} via \texttt{refill\_data}.

Crucially, \texttt{fsm\_mem\_ready} is gated by the interface state
($\texttt{mm\_ready} \wedge \texttt{mi\_state == BUSY}$). Because \texttt{mm\_req}
is de-asserted once accepted, spurious re-triggering is completely prevented
when the FSM transitions from \texttt{ALLOCATE} to \texttt{UPDATE}.

\subsubsection{Top-Level Integration (\texttt{cache\_top.v})}
\texttt{cache\_top} connects all eight sub-modules into a cohesive system,
exposing only CPU-side signals:
\begin{equation}
    \{\texttt{clk}, \texttt{rst}, \texttt{cpu\_req}, \texttt{cpu\_wr}, \texttt{cpu\_addr}, \texttt{cpu\_wdata}\}
    \longrightarrow
    \{\texttt{cpu\_rdata}, \texttt{cpu\_ready}\}
\end{equation}

The cycle-by-cycle operational flows are verified as follows:
\begin{enumerate}
    \item \textbf{Read Hit (1 cycle):} With combinational tag and valid lookup,
        \texttt{hit} asserts in \texttt{IDLE}; \texttt{cpu\_ready} fires in the
        same cycle (fast path).
    \item \textbf{Read Miss (19 cycles total):} \texttt{IDLE} (1) $\to$
        \texttt{COMPARE\_TAG} (1) $\to$ \texttt{ALLOCATE} (14 cycles: 1 interface setup
        + 13 memory burst) $\to$ \texttt{UPDATE} (1 cycle to write data line, tag,
        and valid bit) $\to$ \texttt{RESPOND} (1 cycle asserting \texttt{cpu\_ready})
        $\to$ \texttt{IDLE}.
    \item \textbf{Write Hit (1 cycle):} Fast-path update of the cache word in
        \texttt{data\_ram} while concurrently dispatching the write to memory.
    \item \textbf{Write Miss (16 cycles total):} \texttt{COMPARE\_TAG} $\to$
        \texttt{ALLOCATE} (10 cycles write to memory) $\to$ \texttt{RESPOND}. The
        cache tag and data arrays remain untouched (no-write-allocate).
\end{enumerate}

\subsubsection{Unit and Integration Verification}
Member~3 deliverables are verified via an automated test suite (\texttt{tb/run\_m3\_tests.sh}):
\begin{itemize}
    \item \texttt{tb\_main\_memory.v}: Validates exact 13-cycle read burst,
        10-cycle write, and data integrity across all words.
    \item \texttt{tb\_mem\_interface.v}: Validates request latching, ready pulse
        generation, and absence of phantom re-triggering.
    \item \texttt{tb\_cache\_top.v}: Validates full integration over cold misses,
        spatial locality hits, write-through, no-write-allocate, and conflict-miss
        thrashing (addresses 1\,KB apart).
\end{itemize}
% ---------------------- END PROPOSED SOLUTION C -----------------------


% -------------- REFERENCES (merge into the shared list) ---------------
% \begin{thebibliography}{9}
% \bibitem{smith1982} A.~J. Smith, ``Cache memories,'' \emph{ACM Computing
%   Surveys}, vol.~14, no.~3, pp.~473--530, Sep. 1982.
% \bibitem{hill1989} M.~D. Hill and A.~J. Smith, ``Evaluating associativity in
%   CPU caches,'' \emph{IEEE Trans. Computers}, vol.~38, no.~12,
%   pp.~1612--1630, Dec. 1989.
% \bibitem{jouppi1990} N.~P. Jouppi, ``Improving direct-mapped cache
%   performance by the addition of a small fully-associative cache and
%   prefetch buffers,'' in \emph{Proc. 17th Annu. Int. Symp. Computer
%   Architecture (ISCA)}, 1990, pp.~364--373.
% \bibitem{hennessy2017} J.~L. Hennessy and D.~A. Patterson, \emph{Computer
%   Architecture: A Quantitative Approach}, 6th~ed. Cambridge, MA, USA:
%   Morgan Kaufmann, 2017.
% \bibitem{patterson2020} D.~A. Patterson and J.~L. Hennessy, \emph{Computer
%   Organization and Design RISC-V Edition}, 2nd~ed. Morgan Kaufmann, 2020.
% \end{thebibliography}
 pieces into the IEEEtran report, or copy the
% blocks marked BEGIN/END. Citation keys match m1.tex.
% =====================================================================

% ---------------------- BEGIN LITERATURE REVIEW C ---------------------
\subsection{Write Policies and Main-Memory Interfaces}
\label{sec:litreview-m3}

The interaction between a cache and main memory is governed by two
orthogonal policy choices: \emph{write policy} (when memory is updated) and
\emph{write-allocation policy} (whether a line is brought into cache on a
write miss)~\cite{smith1982, hennessy2017}.

\subsubsection{Write-Through vs.\ Write-Back}
Under a \textbf{write-through} policy, every CPU write immediately updates
both the cache line (if present) and the backing store. The primary
advantage is architectural simplicity and data consistency: main memory is
always up to date, eliminating the need for dirty bits or writebacks on
evictions. Consequently, handling a read-miss refill simply overwrites the
selected line without an eviction writeback~\cite{jouppi1990}. The drawback
is elevated memory-bus traffic, as every store instruction generates a memory
bus transaction regardless of temporal write locality~\cite{patterson2020}.

In contrast, a \textbf{write-back} policy defers memory updates until a
dirty line is evicted from the cache. This significantly reduces write
traffic—particularly for tight inner loops repeatedly modifying local
variables—at the expense of maintaining a per-line \emph{dirty bit} and
implementing a two-phase eviction-refill controller. Phase~II of our project
extends this direct-mapped controller to support write-back.

\subsubsection{Write-Allocate vs.\ No-Write-Allocate}
On a \textbf{write miss}, a \emph{write-allocate} (fetch-on-write) policy
fetches the target block from memory into cache before completing the write,
exploiting spatial and temporal locality under the assumption of subsequent
accesses. Conversely, \emph{no-write-allocate} (write-around) forwards the
write directly to memory, leaving the cache unchanged~\cite{hennessy2017, hill1989}.
In Phase~I, we combine write-through with no-write-allocate: the cleanest and
most deterministic pairing, as writes never trigger block refills or dirty-bit
management.

\subsubsection{Memory Interface Protocols and Burst Latencies}
Modern synchronous DRAM architectures access memory in \emph{bursts}: an
initial row-activation and column-selection latency ($t_{\mathrm{RCD}} + \text{CL}$)
is paid for the first word, after which consecutive words within the open row
are transferred at one word per clock cycle~\cite{hennessy2017}. Our behavioural
memory model parameterises this asymmetry via \texttt{MEM\_FIRST\_LAT} (10 cycles)
and \texttt{MEM\_NEXT\_LAT} (1 cycle). A 4-word (16-byte) block refill thus requires
$10 + 3 \times 1 = 13$ memory cycles, forming the baseline miss penalty for the
AMAT calculations detailed in Section~\ref{sec:analysis-m5}.
% ----------------------- END LITERATURE REVIEW C ----------------------


% --------------------- BEGIN PROPOSED SOLUTION C ----------------------
\subsection{Memory Interface Bridge and Top-Level Integration}
\label{sec:proposed-m3}

\subsubsection{Main-Memory Behavioural Model (\texttt{main\_memory.v})}
The main memory is implemented as a parameterised behavioural model with a
flat 64\,K-word (256\,KB) register array pre-initialised with predictable
word-aligned data ($M[k] = 4k$). The module implements a 3-state controller
(\texttt{IDLE} $\to$ \texttt{FIRST} $\to$ \texttt{BURST}):
\begin{itemize}
    \item \textbf{Block Read:} On \texttt{mem\_req} with \texttt{mem\_wr = 0},
        it latches the block-aligned address and counts \texttt{MEM\_FIRST\_LAT}
        cycles for word~0, followed by \texttt{MEM\_NEXT\_LAT} cycles each for
        words 1 through 3. When all 4 words are assembled, \texttt{mem_ready}
        pulses for 1 cycle with the full 128-bit block on \texttt{mem\_rdata}.
    \item \textbf{Single-Word Write:} On \texttt{mem\_wr = 1}, it waits
        \texttt{MEM\_FIRST\_LAT} cycles, commits the 32-bit word to the array,
        and pulses \texttt{mem\_ready}.
\end{itemize}

\subsubsection{Memory Interface Bridge (\texttt{mem\_interface.v})}
The memory interface decouples the controller FSM from main memory timing. It
latches outgoing requests at the \texttt{IDLE} $\to$ \texttt{BUSY} transition,
asserts \texttt{mm\_req} for one cycle to initiate memory access, holds address
and write-data stable, and routes the returned 128-bit line combinationally
to \texttt{data\_ram.wdata\_line} via \texttt{refill\_data}.

Crucially, \texttt{fsm\_mem\_ready} is gated by the interface state
($\texttt{mm\_ready} \wedge \texttt{mi\_state == BUSY}$). Because \texttt{mm\_req}
is de-asserted once accepted, spurious re-triggering is completely prevented
when the FSM transitions from \texttt{ALLOCATE} to \texttt{UPDATE}.

\subsubsection{Top-Level Integration (\texttt{cache\_top.v})}
\texttt{cache\_top} connects all eight sub-modules into a cohesive system,
exposing only CPU-side signals:
\begin{equation}
    \{\texttt{clk}, \texttt{rst}, \texttt{cpu\_req}, \texttt{cpu\_wr}, \texttt{cpu\_addr}, \texttt{cpu\_wdata}\}
    \longrightarrow
    \{\texttt{cpu\_rdata}, \texttt{cpu\_ready}\}
\end{equation}

The cycle-by-cycle operational flows are verified as follows:
\begin{enumerate}
    \item \textbf{Read Hit (1 cycle):} With combinational tag and valid lookup,
        \texttt{hit} asserts in \texttt{IDLE}; \texttt{cpu\_ready} fires in the
        same cycle (fast path).
    \item \textbf{Read Miss (19 cycles total):} \texttt{IDLE} (1) $\to$
        \texttt{COMPARE\_TAG} (1) $\to$ \texttt{ALLOCATE} (14 cycles: 1 interface setup
        + 13 memory burst) $\to$ \texttt{UPDATE} (1 cycle to write data line, tag,
        and valid bit) $\to$ \texttt{RESPOND} (1 cycle asserting \texttt{cpu\_ready})
        $\to$ \texttt{IDLE}.
    \item \textbf{Write Hit (1 cycle):} Fast-path update of the cache word in
        \texttt{data\_ram} while concurrently dispatching the write to memory.
    \item \textbf{Write Miss (16 cycles total):} \texttt{COMPARE\_TAG} $\to$
        \texttt{ALLOCATE} (10 cycles write to memory) $\to$ \texttt{RESPOND}. The
        cache tag and data arrays remain untouched (no-write-allocate).
\end{enumerate}

\subsubsection{Unit and Integration Verification}
Member~3 deliverables are verified via an automated test suite (\texttt{tb/run\_m3\_tests.sh}):
\begin{itemize}
    \item \texttt{tb\_main\_memory.v}: Validates exact 13-cycle read burst,
        10-cycle write, and data integrity across all words.
    \item \texttt{tb\_mem\_interface.v}: Validates request latching, ready pulse
        generation, and absence of phantom re-triggering.
    \item \texttt{tb\_cache\_top.v}: Validates full integration over cold misses,
        spatial locality hits, write-through, no-write-allocate, and conflict-miss
        thrashing (addresses 1\,KB apart).
\end{itemize}
% ---------------------- END PROPOSED SOLUTION C -----------------------


% -------------- REFERENCES (merge into the shared list) ---------------
% \begin{thebibliography}{9}
% \bibitem{smith1982} A.~J. Smith, ``Cache memories,'' \emph{ACM Computing
%   Surveys}, vol.~14, no.~3, pp.~473--530, Sep. 1982.
% \bibitem{hill1989} M.~D. Hill and A.~J. Smith, ``Evaluating associativity in
%   CPU caches,'' \emph{IEEE Trans. Computers}, vol.~38, no.~12,
%   pp.~1612--1630, Dec. 1989.
% \bibitem{jouppi1990} N.~P. Jouppi, ``Improving direct-mapped cache
%   performance by the addition of a small fully-associative cache and
%   prefetch buffers,'' in \emph{Proc. 17th Annu. Int. Symp. Computer
%   Architecture (ISCA)}, 1990, pp.~364--373.
% \bibitem{hennessy2017} J.~L. Hennessy and D.~A. Patterson, \emph{Computer
%   Architecture: A Quantitative Approach}, 6th~ed. Cambridge, MA, USA:
%   Morgan Kaufmann, 2017.
% \bibitem{patterson2020} D.~A. Patterson and J.~L. Hennessy, \emph{Computer
%   Organization and Design RISC-V Edition}, 2nd~ed. Morgan Kaufmann, 2020.
% \end{thebibliography}
